\documentclass[12pt,oneside]{book}

% Packages for cyberpunk noir styling
\usepackage[utf8]{inputenc}
\usepackage[T1]{fontenc}
\usepackage{geometry}
\usepackage{fancyhdr}
\usepackage{titlesec}
\usepackage{color}
\usepackage{xcolor}
\usepackage{graphicx}
\usepackage{tcolorbox}
\usepackage{enumitem}
\usepackage{hyperref}
\usepackage{fancyvrb}

% Cyberpunk color scheme
\definecolor{neonblue}{RGB}{0,255,255}
\definecolor{neongreen}{RGB}{57,255,20}
\definecolor{darkbg}{RGB}{20,20,20}
\definecolor{matrixgreen}{RGB}{0,200,0}
\definecolor{terminalgreen}{RGB}{0,255,0}

% Page geometry
\geometry{letterpaper,margin=1in}

% Headers and footers
\pagestyle{fancy}
\fancyhf{}
\fancyhead[C]{\color{neonblue}Shadows in the Chain}
\fancyfoot[C]{\color{matrixgreen}\thepage}

% Chapter styling
\titleformat{\chapter}[display]
  {\normalfont\huge\bfseries\color{neongreen}}
  {\chaptertitlename\ \thechapter}
  {20pt}
  {\Huge\color{neonblue}}

% Section styling
\titleformat{\section}
  {\normalfont\Large\bfseries\color{neongreen}}
  {\thesection}
  {1em}
  {}

% Terminal text box styling
\newtcolorbox{terminalbox}{
  colback=black,
  colframe=terminalgreen,
  fontupper=\ttfamily\color{terminalgreen},
  arc=0pt,
  boxrule=1pt
}

% Technical note box
\newtcolorbox{technote}{
  colback=darkbg,
  colframe=neonblue,
  fontupper=\small\color{white},
  arc=3pt,
  boxrule=1pt,
  title=\color{neongreen}\textbf{Technical Note}
}

% Hyperlink setup
\hypersetup{
    colorlinks=true,
    linkcolor=neonblue,
    filecolor=neongreen,
    urlcolor=neonblue,
    pdftitle={Shadows in the Chain - Chapters 3-5},
    pdfauthor={Q-NarwhalKnight Authors},
    pdfsubject={Cyberpunk Noir Novel},
    pdfkeywords={cyberpunk, noir, quantum computing, thriller, cryptography}
}

\begin{document}

\begin{titlepage}
\centering
\vspace*{2cm}

{\Huge\bfseries\color{neongreen} Shadows in the Chain \par}
\vspace{0.5cm}
{\LARGE\color{neonblue} Chapters 3-5 \par}
\vspace{1.5cm}

{\Large\color{neonblue} Q-NarwhalKnight Authors \par}
\vspace{2cm}

{\color{darkbg}
\begin{tcolorbox}[colback=darkbg,colframe=neonblue,arc=0pt,outer arc=0pt]
\color{neongreen}\textbf{Genre:} Cyberpunk, Noir, Thriller\\
\color{neongreen}\textbf{Setting:} Zurich, Singapore, Moscow\\
\color{neongreen}\textbf{Status:} Complete Trilogy\\
\color{neongreen}\textbf{Word Count:} 15,800+ words (Chapters 3-5)
\end{tcolorbox}
}

\vspace{1cm}

{\large\color{matrixgreen}\textbf{Chapter Summary:}\\
$\bullet$ Chapter 3: Quantum Entanglement (Zurich)\\
$\bullet$ Chapter 4: The Singapore Infiltration\\
$\bullet$ Chapter 5: The Moscow Gambit\\[0.5cm]
\textbf{Key Themes:}\\
$\bullet$ Post-quantum cryptography and quantum supremacy\\
$\bullet$ Algorithmic governance vs. human freedom\\
$\bullet$ Impossible choices and moral complexity\\
$\bullet$ Mother-daughter legacy and generational trauma
\par}

\vfill
{\large\color{matrixgreen} Generated by shadowchain-writer CLI + Claude Code \par}
{\color{neonblue} Compiled: 2025-10-09 \par}

\end{titlepage}

\newpage
\tableofcontents
\newpage

% CHAPTER 3 CONTENT
\input{chapter3_content}

\newpage

% CHAPTER 4 CONTENT
\input{chapter4_content}

\newpage

% CHAPTER 5 CONTENT
\input{chapter5_content}

\newpage
\appendix

\chapter{Technical Appendices}

\section{Quantum Computing Threat Timeline}

\begin{technote}
\textbf{Author's Note on Narrative Compression}

This novel accelerates the quantum computing timeline for thriller pacing. Real-world estimates differ significantly from the story's 6-month timeline.
\end{technote}

\subsection{Real-World Timeline (As of 2024)}

\textbf{Current State:}
\begin{itemize}
\item Largest quantum systems: \textasciitilde1,000 physical qubits (IBM, Google, IonQ)
\item Error rates: 0.1\% - 1\% per gate operation
\item Coherence times: microseconds to milliseconds
\item Fault tolerance: Not yet achieved at scale
\end{itemize}

\textbf{Cryptographically Relevant Quantum Computers (CRQC):}
\begin{itemize}
\item Estimated timeline: \textbf{10-15 years} (per NIST, NSA assessments)
\item Required: \textasciitilde20 million physical qubits for RSA-2048
\item Requires: Fault-tolerant error correction at scale
\item Technical barriers: Qubit connectivity, coherence, scalability
\end{itemize}

\textbf{NIST Post-Quantum Standards (Completed 2024):}
\begin{itemize}
\item Crystals-Kyber (key encapsulation)
\item Crystals-Dilithium (digital signatures)
\item SPHINCS+ (stateless signatures)
\item Migration timeline: 5-10 years for full deployment
\end{itemize}

\subsection{Novel's Accelerated Timeline}

\textbf{Kronos Achievement in Story:}
\begin{itemize}
\item 68 logical qubits (fault-tolerant)
\item Scaling to 4,096 qubits in 6 months
\item RSA-2048 breakable in hours (Shor's algorithm)
\item Complete cryptographic supremacy achieved
\end{itemize}

\subsection{Justification for Compression}

\textbf{1. Secret Breakthroughs:}

Kronos has achieved error correction and qubit stability beyond public knowledge. This mirrors classified military/intelligence capabilities being years ahead of public research.

\textbf{2. Classified Programs:}

Nation-state quantum programs (China, US, EU) may be further ahead than public disclosures suggest. Intelligence agencies rarely announce breakthroughs in cryptanalysis.

\textbf{3. Exponential Scaling:}

Quantum computing may follow trajectories similar to Moore's Law once error correction is solved. Early progress is slow; breakthroughs accelerate rapidly.

\textbf{4. Historical Precedent:}

The Manhattan Project achieved nuclear fission 3-5 years faster than mainstream scientific consensus predicted. When resources, talent, and urgency align, technological timelines compress dramatically.

\subsection{Takeaway}

\textcolor{neongreen}{\textbf{Treat the 6-month timeline as ``thriller compression'' of a real 5-10 year threat.}}

The technology is real. The threat is real. The timeline is dramatized for narrative urgency, but the underlying cryptographic crisis is scientifically accurate.

\section{Post-Quantum Cryptography Standards}

\subsection{Crystals-Kyber (Key Encapsulation)}

\textbf{Type:} Lattice-based cryptography

\textbf{Security Level:} Kyber-1024 provides 256-bit quantum security

\textbf{Key Sizes:}
\begin{itemize}
\item Public key: 1,568 bytes
\item Secret key: 3,168 bytes
\item Ciphertext: 1,568 bytes
\end{itemize}

\textbf{Use Case:} Establishing shared secrets for symmetric encryption (replacing RSA key exchange)

\subsection{Crystals-Dilithium (Digital Signatures)}

\textbf{Type:} Lattice-based cryptography

\textbf{Security Level:} Dilithium5 provides 256-bit quantum security

\textbf{Key Sizes:}
\begin{itemize}
\item Public key: 2,592 bytes
\item Secret key: 4,864 bytes
\item Signature: 4,595 bytes
\end{itemize}

\textbf{Use Case:} Digital signatures (replacing RSA, DSA, ECDSA)

\subsection{Migration Strategy}

\textbf{Hybrid Approach (Recommended):}
\begin{itemize}
\item Use both classical and post-quantum algorithms
\item Provides security if either system is broken
\item Gradual transition as PQC matures
\item Maintains backward compatibility
\end{itemize}

\section{Shor's Algorithm}

\textbf{Invented:} 1994 by Peter Shor

\textbf{Impact:} Proves that quantum computers can efficiently factor large integers and solve discrete logarithms---breaking RSA, DSA, and ECDSA.

\subsection{Classical Factorization}

\textbf{Problem:} Given $N = p \times q$ (product of two large primes), find $p$ and $q$.

\textbf{Best Classical Algorithm:} General Number Field Sieve (GNFS)

\textbf{Complexity:} $O(e^{(64/9)^{1/3} (\log N)^{1/3} (\log \log N)^{2/3}})$

For RSA-2048, this is approximately $2^{112}$ operations---infeasible for classical computers.

\subsection{Shor's Algorithm}

\textbf{Quantum Complexity:} $O((\log N)^3)$

This is \textit{polynomial time}---dramatically faster than any known classical algorithm.

\textbf{Requirements for RSA-2048:}
\begin{itemize}
\item Approximately 20 million physical qubits
\item Assuming 0.01\% error rates
\item With surface code error correction
\item Running for several hours
\end{itemize}

\textbf{Kronos's Advantage (In Story):}

With 4,096 logical qubits and aggressive error correction, Kronos has achieved practical RSA-2048 breaking capability. This is well beyond current real-world systems but represents a plausible 10-15 year projection.

\section{Harvest Now, Decrypt Later (HNDL)}

\textbf{Threat Model:}

Adversaries collect encrypted data today, store it, and wait for quantum computers capable of breaking the encryption. Once quantum advantage is achieved, they retroactively decrypt all historical data.

\textbf{Targeted Data:}
\begin{itemize}
\item Government classified communications
\item Corporate trade secrets
\item Personal medical records
\item Financial transactions
\item Long-term strategic intelligence
\end{itemize}

\textbf{Timeline Risk:}

Data encrypted today with RSA-2048 may remain sensitive for 10-20 years. If quantum computers achieve CRQC capability within that timeframe, the data is compromised.

\textbf{Kronos's Strategy (In Story):}

Kronos has been harvesting encrypted data for 3+ years, waiting for quantum supremacy. Now that they've achieved it, they possess decades of retroactively decrypted intelligence---a strategic advantage of unprecedented scale.

\textbf{Mitigation:}

Migrate to post-quantum cryptography \textit{before} CRQC is achieved. This is why the 2024 NIST standards are critical---they provide a head start on the quantum threat.

\chapter{Glossary}

\textbf{BB84:} Quantum key distribution protocol using polarized photons

\textbf{BFT:} Byzantine Fault Tolerance---consensus despite malicious nodes

\textbf{Crystals-Kyber:} NIST post-quantum key encapsulation mechanism

\textbf{Crystals-Dilithium:} NIST post-quantum digital signature scheme

\textbf{CRQC:} Cryptographically Relevant Quantum Computer

\textbf{DAG:} Directed Acyclic Graph---blockchain structure

\textbf{Dilithium5:} Highest security level of Crystals-Dilithium (256-bit)

\textbf{ECDSA:} Elliptic Curve Digital Signature Algorithm (vulnerable to quantum)

\textbf{FSB:} Federal Security Service (Russian intelligence)

\textbf{GCHQ:} Government Communications Headquarters (UK intelligence)

\textbf{HNDL:} Harvest Now, Decrypt Later threat model

\textbf{Kyber-1024:} Highest security level of Crystals-Kyber (256-bit)

\textbf{Lattice-based cryptography:} Post-quantum cryptographic family

\textbf{Logical qubit:} Error-corrected qubit (requires many physical qubits)

\textbf{MSS:} Ministry of State Security (Chinese intelligence)

\textbf{NIST:} National Institute of Standards and Technology

\textbf{NSA:} National Security Agency (US signals intelligence)

\textbf{Physical qubit:} Raw quantum bit (error-prone)

\textbf{PQC:} Post-Quantum Cryptography

\textbf{Q-Day:} Hypothetical date when quantum computers break current cryptography

\textbf{Qubit:} Quantum bit---fundamental unit of quantum information

\textbf{RSA-2048:} 2048-bit RSA encryption (current standard, quantum-vulnerable)

\textbf{Shamir's Secret Sharing:} Threshold cryptographic scheme

\textbf{Shor's Algorithm:} Quantum algorithm for factoring integers

\textbf{SPHINCS+:} NIST post-quantum stateless signature scheme

\textbf{SWIFT:} Society for Worldwide Interbank Financial Telecommunication

\end{document}
